% !TeX root = ../main.tex

\chapter{Introduction}\label{chapter:introduction}

Modern high-performance and parallel applications increasingly face dynamic and heterogeneous execution environments, where the optimal allocation of computational resources may change over time. While OpenMP~\parencite{dagum_openmp_1998} provides a high-level programming model for shared-memory parallelism, its execution model largely assumes static or monotonically increasing resource allocation within a process.

Although the OpenMP standard has included the \texttt{OMP\_DYNAMIC} mechanism for more than two decades~\parencite{openmp60}, dynamic adjustment of computational resources has not been fully realized in practice. In particular, current OpenMP runtimes typically rely on fixed or lazily growing thread pools, without providing standardized mechanisms to dynamically adapt the number of active threads across parallel regions or in response to external resource constraints.

The goal of this thesis is to investigate and prototype dynamic resource management for OpenMP, with a primary focus on CPU cores and an outlook towards accelerator resources such as GPUs. The work is based on an in-depth analysis of the LLVM OpenMP runtime (libomp)~\parencite{noauthor_welcome_nodate}, including the lifecycle of worker threads, thread pool behavior, and relevant control points within the runtime.

Building on this analysis, different approaches to dynamic resource management are explored. This includes the integration of an external resource manager via a communication mechanism as a proof of concept, as well as the design and implementation of a custom, runtime-internal resource manager supporting multiple allocation strategies. The feasibility, overhead, and limitations of these approaches are evaluated with respect to the OpenMP execution model and standard compliance.

Finally, the proposed mechanisms are evaluated using microbenchmarks and application benchmarks to assess performance impact, overhead, and compatibility. Based on the experimental findings, this thesis discusses implications for OpenMP runtimes and derives a proposal for standardizing dynamic resource management interfaces in future versions of the OpenMP standard.
